\documentclass[11pt]{article}
\usepackage{amsmath,amssymb,amsthm}
\newtheorem{theorem}{Theorem}
\newtheorem{lemma}{Lemma}
\newtheorem{corollary}{Corollary}
\newtheorem{remark}{Remark}
\newcommand{\EMD}{\mathrm{EMD}}

\title{Positivity of $h_m$ for $3\mid d$, and an Orbit--Product Formula for $h_m$ when $\gcd(d,3)=1$}
\author{Seed 2, Layer 2, Round 2}
\date{2026-07-04}

\begin{document}
\maketitle

\begin{abstract}
Let $c=(c_0,c_1,c_2)$ be a profile with $d=c_0+c_1+c_2$, let $F_{c,m}(q)$ be the generating
function of cylindric partitions of profile $c$ with maximum entry at most $m$, let
$g_m = F_{c,m}-F_{c,m-1}$ and $h_m = (q^\ell;q^\ell)_m\, g_m$ with $\ell=\gcd(d,3)$.
We prove two results. (1) For $3\mid d$, $h_m$ has nonnegative coefficients; this settles the
$\ell=3$ half of the core bottleneck identified in the Layer-1 synthesis. (2) For
$\gcd(d,3)=1$ we prove an exact \emph{orbit--product formula}: $h_m$ (and also
$P_m := (q;q)_m F_{c,m}$) is a finite sum, over sequences of $C_3$-orbits of profiles, of
products $\prod_{j=1}^m U^{(j)}(q^j)$ of explicit integer polynomials $U^{(j)}$ with
coefficients in $\{0,\pm 1\}$. In particular $(q;q)_m F_{c,m} \in \mathbb{Z}[q]$, and
$h_m(1)=K^m$ with $K = \binom{d+2}{2}/3$. The remaining obstacle to $h_m\ge 0$ for
$\gcd(d,3)=1$ is isolated as the positivity of this signed orbit-path sum.
Both results take as input only the Adjugate Monomial Theorem
$\operatorname{adj}(I-A(x))[c,c'] = x^{\EMD(c,c')}$ and $\det(I-A(x)) = -(x^3-1)$,
which were proved in Layer 1 (Seed 4).
\end{abstract}

\section{Setup and inputs}

Profiles of weight $d$ are triples $c=(c_0,c_1,c_2)\in\mathbb{Z}_{\ge0}^3$ with $\sum c_i = d$.
The Corteel--Welsh recursion for the bounded generating functions $F_{c,m}$ is encoded by the
matrix $A(x)$ with entries indexed by profiles, and Layer 1 (Seed 4) proved:

\medskip
\noindent\textbf{Input 1 (Adjugate Monomial Theorem, GREEN).}
$\operatorname{adj}(I-A(x))[c,c'] = x^{\EMD(c,c')}$, where
\begin{equation}\label{eq:emd}
\EMD(c,c') \;=\; 3\max\bigl(0,\; c'_1-c_1,\; c_0-c'_0\bigr) + (c'_0-c_0) - (c'_1-c_1),
\end{equation}
and $\det(I-A(x)) = -(x^3-1)$. Consequently, for every $m\ge 1$,
\begin{equation}\label{eq:inversion}
(1-q^{3m})\, F_{c,m} \;=\; \sum_{c'} q^{m\,\EMD(c,c')}\, F_{c',m-1},
\qquad F_{c,0}=1 .
\end{equation}

\noindent $\EMD$ is the minimum-cost clockwise transport distance on $\mathbb{Z}/3$:
writing $f=(f_0,f_1,f_2)$ for a clockwise flow ($f_i$ units crossing edge $i\to i+1$),
the transport equations are $c_i - c'_i = f_i - f_{i-1}$, whose solutions form the line
$f = f^\ast + t(1,1,1)$; then $\EMD(c,c') = \min\{\,f_0+f_1+f_2 : f\ge 0\,\}$, and
\eqref{eq:emd} is the closed form. Note $\EMD(c,c)=0$.

\section{Symmetry lemmas}

Let $\sigma$ denote the cyclic rotation $\sigma(c_0,c_1,c_2)=(c_2,c_0,c_1)$.

\begin{lemma}[Profile-rotation invariance of $F$]\label{lem:Finv}
$F_{\sigma c, m} = F_{c,m}$ for all $c$ and all $m\ge 0$ (as power series in $q$).
\end{lemma}

\begin{proof}
A cylindric partition of profile $(c_1,\dots,c_k)$ is a tuple $(\lambda^{(1)},\dots,\lambda^{(k)})$
with $\lambda^{(i)}_j \ge \lambda^{(i+1)}_{j+c_{i+1}}$ ($1\le i\le k-1$) and
$\lambda^{(k)}_j \ge \lambda^{(1)}_{j+c_1}$. The relabelling
$\Lambda=(\lambda^{(1)},\dots,\lambda^{(k)})\mapsto
\Lambda'=(\lambda^{(2)},\dots,\lambda^{(k)},\lambda^{(1)})$
is a bijection from cylindric partitions of profile $(c_1,\dots,c_k)$ to those of profile
$(c_2,\dots,c_k,c_1)$: each defining inequality of $\Lambda'$ is one of the defining
inequalities of $\Lambda$ (the interior ones shift the index by one; the two boundary
inequalities swap roles). The bijection preserves total size and maximum entry, hence
preserves $F_{\cdot,m}$. Iterating gives invariance under the full rotation group $C_3$,
in particular under $\sigma$.
\end{proof}

\begin{lemma}[Rotation invariance of $\EMD$]\label{lem:emdinv}
$\EMD(\sigma c, \sigma c') = \EMD(c,c')$.
\end{lemma}

\begin{proof}
Let $a = c'_0-c_0$, $b = c'_1-c_1$, so $\EMD(c,c') = 3\max(0,b,-a)+a-b$ by \eqref{eq:emd}.
For the rotated pair, $(\sigma c')_0-(\sigma c)_0 = c'_2-c_2 = -(a+b)$ (since the coordinate
differences sum to $0$) and $(\sigma c')_1-(\sigma c)_1 = c'_0-c_0 = a$. Hence
\[
\EMD(\sigma c,\sigma c') = 3\max\bigl(0,\, a,\, a+b\bigr) - (a+b) - a
= 3\bigl(a + \max(-a, 0, b)\bigr) - 2a - b
= 3\max(0,b,-a) + a - b,
\]
which is $\EMD(c,c')$.
\end{proof}

\begin{lemma}[Mod-3 shift]\label{lem:mod3}
$\EMD(c,\sigma c') \equiv \EMD(c,c') + d \pmod 3$. Consequently, if $\gcd(d,3)=1$ the three
values $\EMD(c,c'), \EMD(c,\sigma c'), \EMD(c,\sigma^2 c')$ are pairwise distinct mod $3$
and represent all residues $\{0,1,2\}$; if $3\mid d$ they are all congruent mod $3$.
\end{lemma}

\begin{proof}
Modulo 3, \eqref{eq:emd} gives $\EMD(c,c') \equiv (c'_0-c_0)-(c'_1-c_1)$. Hence
\[
\EMD(c,\sigma c') - \EMD(c,c') \equiv \bigl[(c'_2-c_0)-(c'_0-c_1)\bigr]
- \bigl[(c'_0-c_0)-(c'_1-c_1)\bigr]
= c'_1+c'_2-2c'_0 \equiv d - 3c'_0 \equiv d \pmod 3. \qedhere
\]
\end{proof}

(Lemmas \ref{lem:emdinv} and \ref{lem:mod3} were additionally verified numerically for all
pairs of profiles with $d\le 20$.)

\section{Theorem 1: $h_m \ge 0$ for $3 \mid d$}

Define, for any $d$, $\widetilde P_m(c) := (q^3;q^3)_m\, F_{c,m}$.

\begin{lemma}\label{lem:Ptilde}
For all $d$ and all $m \ge 0$, $\widetilde P_m(c)$ has nonnegative coefficients; indeed
\begin{equation}\label{eq:Ptilderec}
\widetilde P_m(c) = \sum_{c'} q^{m\,\EMD(c,c')}\, \widetilde P_{m-1}(c'), \qquad \widetilde P_0 = 1 .
\end{equation}
\end{lemma}

\begin{proof}
Multiply \eqref{eq:inversion} by $(q^3;q^3)_{m-1}$ and use
$(q^3;q^3)_m = (1-q^{3m})(q^3;q^3)_{m-1}$. Nonnegativity follows by induction on $m$,
since \eqref{eq:Ptilderec} has nonnegative coefficients on the right-hand side.
(Unfolding \eqref{eq:Ptilderec} recovers the manifestly positive EMD path formula of
Round 1, Agent B.)
\end{proof}

\begin{theorem}\label{thm:div3}
Let $3\mid d$, so $\ell=\gcd(d,3)=3$ and $h_m = (q^3;q^3)_m\,(F_{c,m}-F_{c,m-1})$. Then for
all profiles $c$ and all $m\ge 1$, $h_m$ has nonnegative coefficients. Explicitly,
\begin{equation}\label{eq:hdiv3}
h_m(c) \;=\; \sum_{c'\ne c} q^{m\,\EMD(c,c')}\, \widetilde P_{m-1}(c') \;+\; q^{3m}\, \widetilde P_{m-1}(c).
\end{equation}
\end{theorem}

\begin{proof}
By definition and Lemma~\ref{lem:Ptilde},
\[
h_m = (q^3;q^3)_m F_{c,m} - (q^3;q^3)_m F_{c,m-1}
= \widetilde P_m(c) - (1-q^{3m})\,\widetilde P_{m-1}(c).
\]
Substitute \eqref{eq:Ptilderec} for $\widetilde P_m(c)$; the $c'=c$ term equals
$q^{m\,\EMD(c,c)}\widetilde P_{m-1}(c) = \widetilde P_{m-1}(c)$, and
$\widetilde P_{m-1}(c) - (1-q^{3m})\widetilde P_{m-1}(c) = q^{3m}\widetilde P_{m-1}(c)$,
giving \eqref{eq:hdiv3}. Every term on the right of \eqref{eq:hdiv3} is a nonnegative
series by Lemma~\ref{lem:Ptilde}.
\end{proof}

\begin{remark}
Theorem~\ref{thm:div3} upgrades the $3\mid d$ half of the Layer-1 core bottleneck
(``$(q^\ell;q^\ell)_m g_m \ge 0$ for all $d$'') from YELLOW to GREEN. It is a strictly
one-line consequence of the Adjugate Monomial Theorem; the difficulty of the conjectural
case $\gcd(d,3)=1$ is therefore *entirely* concentrated in the division by
$\prod_{j=1}^m (1+q^j+q^{2j}) = (q^3;q^3)_m/(q;q)_m$.
\end{remark}

\section{Theorem 2: the orbit--product formula for $\gcd(d,3)=1$}

Assume now $\gcd(d,3)=1$ and set $P_m(c) := (q;q)_m F_{c,m}$, so $h_m = P_m - (1-q^m)P_{m-1}$.
Since $d\not\equiv 0 \pmod 3$, no profile satisfies $\sigma c = c$, so $\sigma$ acts freely
and the $\binom{d+2}{2}$ profiles fall into $K := \binom{d+2}{2}/3$ orbits
$O = \{c',\sigma c',\sigma^2 c'\}$.

\begin{lemma}[Local divisibility]\label{lem:localdiv}
For profiles $c, c'$ define
$T_{c,[c']}(x) := \sum_{r=0}^{2} x^{\EMD(c,\sigma^r c')}$, and the top-modified version
$T^{\mathrm{top}}_{c,[c']}(x)$, which is defined as $T$ except that when $c \in [c']$
the term $x^{\EMD(c,c)} = x^0$ is replaced by $x^3$. Then $T$ and $T^{\mathrm{top}}$ depend
only on $c$ and the orbit $[c']$, and both are divisible by $1+x+x^2$ in $\mathbb{Z}[x]$.
We write $U_{c,[c']} := T_{c,[c']}/(1+x+x^2)$ and
$U^{\mathrm{top}}_{c,[c']} := T^{\mathrm{top}}_{c,[c']}/(1+x+x^2)$.
\end{lemma}

\begin{proof}
Independence of the representative $c'$ is clear ($\sigma^r c'$ runs over the orbit).
By Lemma~\ref{lem:mod3} the three exponents represent all residues mod 3.
Since $x^k \equiv x^{k \bmod 3} \pmod{1+x+x^2}$ and $1+x+x^2 \equiv 0$, any sum of three
monomials whose exponents cover all residues mod 3 is divisible by $1+x+x^2$. The top
modification replaces the exponent $0$ by $3$, preserving the residue, so divisibility
persists. (Independence of representative in the top case: the modified term is attached to
the unique representative equal to $c$ itself.)
\end{proof}

\begin{theorem}[Orbit--product formula]\label{thm:orbit}
Let $\gcd(d,3)=1$ and let $\mathcal{O}$ be the set of $\sigma$-orbits of profiles of weight
$d$. Then for all $m\ge 0$ and all profiles $c$,
\begin{equation}\label{eq:orbitP}
P_m(c) \;=\; \sum_{(O_{m-1},\dots,O_0)\,\in\,\mathcal{O}^m}\;
\prod_{j=1}^{m} U_{\,o_j,\,O_{j-1}}(q^j),
\end{equation}
where $o_m := c$ and $o_j$ denotes any representative of $O_j$ for $j<m$ (the product is
well defined by Lemma~\ref{lem:localdiv} and Lemma~\ref{lem:emdinv}); and
\begin{equation}\label{eq:orbith}
h_m(c) \;=\; \sum_{(O_{m-1},\dots,O_0)\,\in\,\mathcal{O}^m}\;
U^{\mathrm{top}}_{\,c,\,O_{m-1}}(q^m)\,
\prod_{j=1}^{m-1} U_{\,o_j,\,O_{j-1}}(q^j).
\end{equation}
\end{theorem}

\begin{proof}
First, one level. By Lemma~\ref{lem:Finv}, $P_{m-1}(\sigma c')=P_{m-1}(c')$, so grouping the
right side of \eqref{eq:inversion} (multiplied by $(q;q)_{m-1}$) into orbits,
\[
(1-q^{3m})\,(q;q)_{m-1} F_{c,m} \;=\; \sum_{O\in\mathcal{O}} T_{c,O}(q^m)\, P_{m-1}(O).
\]
The left side equals $(1+q^m+q^{2m})\,P_m(c)$, and by Lemma~\ref{lem:localdiv}
$T_{c,O}(q^m) = (1+q^m+q^{2m})\,U_{c,O}(q^m)$. Since $\mathbb{Z}[[q]]$ is an integral domain
and $1+q^m+q^{2m}\neq 0$, we may cancel:
\begin{equation}\label{eq:onelevel}
P_m(c) \;=\; \sum_{O\in\mathcal{O}} U_{c,O}(q^m)\, P_{m-1}(O).
\end{equation}
Well-definedness in $O$ of the coefficient uses Lemma~\ref{lem:emdinv}: replacing the
representative of an intermediate orbit rotates both arguments of the relevant $\EMD$'s
simultaneously and changes nothing. Iterating \eqref{eq:onelevel} down to $P_0=1$ gives
\eqref{eq:orbitP}.

For \eqref{eq:orbith}: $h_m = P_m(c) - (1-q^m)P_{m-1}(c)$. Multiply by $1+q^m+q^{2m}$:
\[
(1+q^m+q^{2m})\,h_m = \sum_{O} T_{c,O}(q^m) P_{m-1}(O) - (1-q^{3m}) P_{m-1}(c).
\]
In the orbit $O_c$ of $c$, the term $q^{\EMD(c,c)\,m}P_{m-1}(c) = P_{m-1}(c)$ combines with
$-(1-q^{3m})P_{m-1}(c)$ to give $q^{3m}P_{m-1}(c)$: this is exactly the replacement
$x^0 \mapsto x^3$ defining $T^{\mathrm{top}}$. Hence
\[
(1+q^m+q^{2m})\,h_m = \sum_{O} T^{\mathrm{top}}_{c,O}(q^m)\, P_{m-1}(O),
\]
and cancelling $1+q^m+q^{2m}$ (Lemma~\ref{lem:localdiv}) and expanding $P_{m-1}$ by
\eqref{eq:orbitP} gives \eqref{eq:orbith}.
\end{proof}

\begin{corollary}\label{cor:poly}
For $\gcd(d,3)=1$: $P_m = (q;q)_m F_{c,m}$ and $h_m$ lie in $\mathbb{Z}[q]$ (they are honest
polynomials), and $P_m(1) = h_m(1) = K^m$ where $K=\binom{d+2}{2}/3 = (d+1)(d+2)/6$.
\end{corollary}

\begin{proof}
Each $U$, $U^{\mathrm{top}}$ is an integer polynomial and the sums in
Theorem~\ref{thm:orbit} are finite. At $q=1$: $T_{c,O}(1)=3$ and $T^{\mathrm{top}}_{c,O}(1)=3$,
so every $U(1) = U^{\mathrm{top}}(1) = 1$, and each of the $K^m$ orbit sequences contributes
exactly $1$.
\end{proof}

\begin{remark}[Computed structure of the $U$'s]
Numerically (all $d\le 14$ with $\gcd(d,3)=1$, all pairs $(c,O)$, both plain and top):
every $U_{c,O}$ has coefficients in $\{0,\pm1\}$, and its nonzero coefficients alternate in
sign, beginning and ending with $+1$ (so $U = x^{a_0}-x^{a_1}+\cdots+x^{a_{2t}}$ with
$a_0<a_1<\cdots<a_{2t}$ and $U(1)=1$). Also verified: \eqref{eq:orbitP} and \eqref{eq:orbith}
reproduce the transfer-matrix values of $P_m$ and $h_m$ exactly for $d\in\{1,2,4,5\}$,
all orbits, $m\le 3$ ($m\le 2$ for $d=5$), and the resulting $h$-values agree with Seed 8's
independently computed tables (e.g.\ $d=4$, $c=(2,1,1)$: $h_2 = [0,0,3,4,5,3,3,2,2,1,1,0,1]$).
\end{remark}

\begin{remark}[What remains]
The single remaining obstacle to the core bottleneck ($h_m\ge0$, hence
$Q_n\ge 0$ via the $D_k^m$ tower) is the positivity of the signed sum \eqref{eq:orbith}.
Three natural strengthenings FAIL and should not be retried na\"ively:
(i) individual orbit-sequence products $\prod_j U(q^j)$ can have negative coefficients;
(ii) Abel summation along a domination chain fails because the orbit values
$P_{m-1}(O)$ are pairwise incomparable coefficientwise for $m\ge 2$;
(iii) shifted domination $q^\delta P(O) \le P(O')$ fails for all $\delta$ (the $P$'s are
polynomials of comparable degree). What is promising: the $U$'s are $\{0,\pm1\}$
alternating with $U(1)=1$, evaluated at distinct scales $q^j$ --- the sum \eqref{eq:orbith}
is a signed lattice-path model begging for a sign-reversing involution that fixes exactly
$K^m$ monomials (one per orbit sequence).
\end{remark}

\end{document}
